% from ega2-1-fr.tex
\oldpage[II]{40}
que si $k>m$, on a $S_kx_i=0$ pour tout $i$~; si $h$ est le plus grand des
degrés des $x_i$, on en conclut que $M_k=0$ pour $k>h+m$, ce qui termine la
démonstration.
\end{proof}

\begin{corollary}[2.7.4]
\label{II.2.7.4-fr}
Soit $S$ un anneau gradué tel que l'idéal $S_+$ soit de type fini~; pour que
$X=\operatorname{Proj}(S)=\varnothing$, il faut et il suffit qu'il existe $n$
tel que $S_k=0$ pour $k\geq n$.
\end{corollary}

\begin{proof}
La condition $X=\varnothing$ est en effet équivalente à
$\mathcal{O}_X=\widetilde{S}=0$, et $S$ est un $S$-module monogène.
\end{proof}

\begin{theorem}[2.7.5]
\label{II.2.7.5-fr}
On suppose que l'idéal $S_+$ est engendré par un nombre fini d'éléments
homogènes de degré $1$~; soit $X=\operatorname{Proj}(S)$. Alors, pour tout
$\mathcal{O}_X$-Module quasi-cohérent $\mathcal{F}$, l'homomorphisme
canonique $\beta:\widetilde{\Gamma_*(\mathcal{F})}
\to\mathcal{F}$ (\hyperref[II.2.6.4-fr]{2.6.4}) est un isomorphisme.
\end{theorem}

\begin{proof}
En effet, si $S_+$ est engendré par un nombre fini d'éléments $f_i\in S_1$,
$X$ est réunion des sous-espaces $\operatorname{Spec}(S_{(f_i)})$
(\hyperref[II.2.3.6-fr]{2.3.6}) qui sont quasi-compacts, donc est
quasi-compact~; en outre $X$ est un schéma
(\hyperref[II.2.4.2-fr]{2.4.2})~; d'après
(\hyperref[I.9.3.2-fr]{I, 9.3.2}),
(\hyperref[II.2.5.14.2-fr]{2.5.14.2}) et
(\hyperref[II.2.6.3-fr]{2.6.3}) on a, pour tout $f\in S_d$ ($d>0$), un
isomorphisme canonique
$(\Gamma_*(\mathcal{F}))_{(\alpha_d(f))}\xrightarrow{\sim}
\Gamma(D_+(f),\mathcal{F})$~; d'ailleurs, par définition,
$(\Gamma_*(\mathcal{F}))_{(\alpha_d(f))}$ (où $\Gamma_*(\mathcal{F})$ est
considéré comme $\Gamma_*(\mathcal{O}_X)$-module) n'est autre que
$(\Gamma_*(\mathcal{F}))_{(f)}$ (où $\Gamma_*(\mathcal{F})$ est considéré
comme $S$-module)~; si on se reporte à la définition
(\hyperref[I.9.3.1-fr]{I, 9.3.1}) de l'isomorphisme canonique précédent, on
voit qu'il coïncide avec l'homomorphisme $\beta_f$, d'où le théorème.
\end{proof}

\begin{remark}[2.7.6]
\label{II.2.7.6-fr}
Si on suppose l'anneau gradué $S$ \emph{noethérien}, la condition de
(\hyperref[II.2.7.5-fr]{2.7.5}) est vérifiée \emph{ipso facto} dès que l'on
suppose l'idéal $S_+$ engendré par l'ensemble $S_1$ des éléments homogènes de
degré $1$.
\end{remark}

\begin{corollary}[2.7.7]
\label{II.2.7.7-fr}
Sous les hypothèses de (\hyperref[II.2.7.5-fr]{2.7.5}), tout
$\mathcal{O}_X$-Module quasi-cohérent $\mathcal{F}$ est isomorphe à un
$\mathcal{O}_X$-Module de la forme $\widetilde{M}$, où $M$ est un $S$-module
gradué.
\end{corollary}

\begin{corollary}[2.7.8]
\label{II.2.7.8-fr}
Sous les hypothèses de (\hyperref[II.2.7.5-fr]{2.7.5}), tout
$\mathcal{O}_X$-Module quasi-cohérent de type fini $\mathcal{F}$ est
isomorphe à un $\mathcal{O}_X$-Module de la forme $\widetilde{N}$, où $N$ est
un $S$-module gradué de type fini.
\end{corollary}

\begin{proof}
On peut supposer $\mathcal{F}=\widetilde{M}$, où $M$ est un $S$-module gradué
(\hyperref[II.2.7.7-fr]{2.7.7}). Soit $(f_\lambda)_{\lambda\in L}$ un
système de générateurs homogènes de $M$~; pour toute partie finie $H$ de $L$,
soit $M_H$ le sous-module gradué de $M$ engendré par les $f_\lambda$ tels que
$\lambda\in H$~; il est clair que $M$ est limite inductive de ses sous-modules
$M_H$, donc $\mathcal{F}$ est limite inductive de ses sous-$\mathcal{O}_X$-Modules
$\widetilde{M_H}$ (\hyperref[II.2.5.4-fr]{2.5.4}). Mais comme $\mathcal{F}$
est de type fini et l'espace sous-jacent $X$ quasi-compact, il résulte de
(\hyperref[0.5.2.3-fr]{0, 5.2.3}) que $\mathcal{F}=\widetilde{M_H}$ pour une
partie finie $H$ de $L$.
\end{proof}

\begin{corollary}[2.7.9]
\label{II.2.7.9-fr}
Sous les hypothèses de (\hyperref[II.2.7.5-fr]{2.7.5}), soit $\mathcal{F}$ un
$\mathcal{O}_X$-Module quasi-cohérent de type fini. Il existe alors un entier
$n_0$ tel que, pour tout $n\geq n_0$, $\mathcal{F}(n)$ soit isomorphe à un
quotient d'un $\mathcal{O}_X$-Module de la forme $\mathcal{O}_X^k$ ($k>0$
dépendant de $n$), et par suite est engendré par un nombre fini de ses sections
au-dessus de $X$ (\hyperref[0.5.1.1-fr]{0, 5.1.1}).
\end{corollary}

\begin{proof}
En vertu de (\hyperref[II.2.7.8-fr]{2.7.8}), on peut supposer que
$\mathcal{F}=\widetilde{M}$, où $M$ est quotient d'une somme directe finie de
$S$-modules de la forme $S(m_i)$~; en vertu de
(\hyperref[II.2.5.4-fr]{2.5.4}) on peut donc se borner au cas où $M=S(m)$,
donc $\mathcal{F}(n)=\widetilde{S(m+n)}
=\mathcal{O}_X(m+n)$ (\hyperref[II.2.5.15-fr]{2.5.15}). Il suffira donc de
démontrer le

